%------------------------------------------------------------------------------%
\documentclass[fleqn,utf8,aspectratio=169,12pt]{beamer}

\usepackage{calculo_varias_variaveis-1}

\title{Propriedades do Vetor Gradiente}

%------------------------------------------------------------------------------%
\begin{document}

\addtitle

%------------------------------------------------------------------------------%
\section{Vetor Gradiente}

%------------------------------------------------------------------------------%
\begin{frame}[t]{Vetor Gradiente}
  Se uma função \(f(x, y)\) é derivável,
  seu \emph{vetor gradiente} é

  \vspace{0.99\baselineskip}
  \[
    \nabla f
    \;=\; \D{f}{x}\veci + \D{f}{y}\vecj
    \;=\; \left(\begin{array}{c}
            \displaystyle\D{f}{x} \\[5mm]
            \displaystyle\D{f}{y}
          \end{array}\right)
  \]
\end{frame}

%------------------------------------------------------------------------------%
\begin{frame}[t]{Vetor Gradiente 3D}
  Se uma função \(f(x, y, z)\) é derivável,
  seu \emph{vetor gradiente} é

  \[
    \nabla f
    \;=\; \D{f}{x}\veci + \D{f}{y}\vecj + \D{f}{z}\veck
    \;=\; \left(\begin{array}{c}
            \displaystyle\D{f}{x} \\[5mm]
            \displaystyle\D{f}{y} \\[5mm]
            \displaystyle\D{f}{z}
          \end{array}\right)
  \]
\end{frame}

%------------------------------------------------------------------------------%
\section{Gradientes e Crescimento da Função}

%------------------------------------------------------------------------------%
\begin{proofframe}{Derivada Direcional}
  \begin{align*}
    \onslide<+->{D_u f}
    \onslide<+->{& = \nabla f \cdot u                        \\[5mm]}
    \onslide<+->{& = \norm{\nabla f}\; \norm{u} \cos(\theta) \\[5mm]}
    \onslide<+->{& = \norm{\nabla f} \cos(\theta)}
  \end{align*}
\end{proofframe}

%------------------------------------------------------------------------------%
\begin{proofframe}{Crescimento da Função}
\begin{enumerate}
  \setlength\itemsep{1.5em}
  \item A função aumenta mais rapidamente quando
        \pause
        \tabto{82mm}\emph{$\cos(\theta)=1$},
        \pause
        \emph{$\hphantom{-}\quad\theta=0$} \\[2mm]
        \pause
        Direção do gradiente
        \pause
  \item A função descresse mais rapidamente quando
        \pause
        \tabto{82mm}\emph{$\cos(\theta)=-1$},
        \pause
        \emph{$\quad\theta=\pi$} \\[2mm]
        \pause
        Direção oposta ao gradiente
        \pause
  \item Se $u$ for \emph{ortogonal} a \(\nabla f\neq 0\)
        \pause
        ele aponta para uma direção de \emph{variação nula}
\end{enumerate}
\end{proofframe}

%------------------------------------------------------------------------------%
\begin{frame}{Gradientes e Crescimento da Função}
  \begin{enumerate}[<+->]
    \setlength\itemsep{1.5em}
    \item Direção de \emph{maior crescimento} da função $f$ no ponto \((x, y)\)
          \[
            v_M = \nabla f(x, y)
          \]
    \item Direção de \emph{maior decrescimento} da função $f$ no ponto \((x, y)\)
          \[
            v_m = -\nabla f(x, y)
          \]
    \item Direção de \emph{variação nula} da função $f$ no ponto \((x, y)\)
          \[
            \nabla f(x, y) \cdot v_o  = 0
          \]
  \end{enumerate}
\end{frame}

%------------------------------------------------------------------------------%
\include{ex/H-02-Propriedades_gradientes-ex1}

%------------------------------------------------------------------------------%
\section{Gradientes e Curvas de Nível}

%------------------------------------------------------------------------------%
\begin{proofframe}{Curvas de Nível}
  Curva de Nível é o conjunto de pontos \((x, y)\) onde
  \[
    f(x, y) = c
  \]
  para um valor constante $c$

  \pause
  \vspace{0.5\baselineskip}
  Em alguns casos, é possível escrever essa região como uma curva paramétrica
  \[
    x = g(t)
    \qquad
    y = h(t)
  \]
  \pause
  Vetorialmente temos
  \[
    r(t)
    = g(t)\veci + h(t)\vecj
    = \vetor{g(t)}{h(t)}
  \]
\end{proofframe}

%------------------------------------------------------------------------------%
\begin{proofframe}{Tangente de uma Curva de Nível}
\begin{columns}
\column{0.5\linewidth}
  \begin{align*}
    \onslide<+->{f\big(g(t), h(t)\big) & = c \\[2mm]}
    \onslide<+->{\frac{d}{dt} f\big(g(t), h(t)\big) & = 0 \\[2mm]}
    \onslide<+->{\D{f}{x}\frac{dg}{dt} + \D{f}{y}\frac{dh}{dt} & = 0  \\[2mm]}
    \onslide<+->{\Vetor{\D{f}{x}}{\D{f}{y}} \cdot
                 \Vetor{\frac{dg}{dt}}{\frac{dh}{dt}} & = 0}
  \end{align*}
\column{0.5\linewidth}
  \onslide<+->{\[\nabla f \cdot \frac{dr}{dt} = 0\]}

  \onslide<+->{O gradiente é \emph{ortogonal} ao vetor tangente à curva de nível}
\end{columns}
\end{proofframe}

%------------------------------------------------------------------------------%
\begin{frame}{Funções Diferenciáveis}
  Em todo ponto \((a, b)\) no domínio de uma \emph{função diferenciável}
  \(f(x, y)\),

  \pause
  o \emph{gradiente} de $f$
  \pause
  é normal à \emph{curva de nível} que passa por \((a, b)\)
\end{frame}

%------------------------------------------------------------------------------%
\framedgraphic{Gradientes e Curvas de Nível}{gradiente_curva_nivel}{}

%------------------------------------------------------------------------------%
%------------------------------------------------------------------------------%
\begin{question}
  Dada a função
  \[
    f(x, y) = 6x - 3y + 9
  \]

  a) Determine e esboce a curva de nível \(f(x, y) = 6\)

  b) Calcule e esboce o gradiente de $f$ nos pontos
  \((0, 1)\) e
  \((1, 3)\)
\end{question}

%------------------------------------------------------------------------------%
\begin{resolution}{Solução}
  \onslide<+->{a) Determinando a curva de nível}
  \begin{align*}
    \onslide<+->{f(x, y)    & = 6}   \\
    \onslide<+->{6x -3y + 9 & = 6} \\
    \onslide<+->{-3y & = 6 -6x - 9} \\
    \onslide<+->{-3y & = -6x -3} \\
    \onslide<+->{  y & =  2x + 1}
  \end{align*}
  \onslide<+->{A curva de nível é uma reta}
\end{resolution}

%------------------------------------------------------------------------------%
\begin{resolution}{Solução}
  b) Calculando as derivadas parciais de $f$

  \pause
  \[
    \D{f}{x}
    \pause
    \;=\; \D{}{x}\left(6x - 3y + 9\right)
    \pause
    \;=\; 6
  \]

  \pause
  \[
    \D{f}{y}
    \pause
    \;=\; \D{}{y}\left(6x - 3y + 9\right)
    \pause
    \;=\; -3
  \]

  \pause
  Gradiente de $f$
  \pause
  \[
    \nabla f(x, y)
    \pause
    \;=\; \nabla\left(6x - 3y + 9\right)
    \pause
    \;=\; \vetor{6}{-3}
  \]
\end{resolution}

%------------------------------------------------------------------------------%
\begin{resolution}{Solução}
  Calculando o gradiente nos pontos
  \((0, 1)\) e
  \((1, 3)\)

  \pause
  \[
    \nabla f(0, 1)
    \pause
    \;=\; \vetor{6}{-3}
  \]

  \pause
  \[
    \nabla f(2, -1)
    \pause
    \;=\; \vetor{6}{-3}
  \]
\end{resolution}

%------------------------------------------------------------------------------%
%------------------------------------------------------------------------------%
\begin{question}
  Dada a função
  \[
    f(x, y) = (x-2)^2 + (y-1)^2
  \]

  a) Determine e esboce a curva de nível \(f(x, y) = 4\)

  b) Calcule e esboce o gradiente de $f$ nos pontos
  \((0,  1)\),
  \((2, -1)\),
  \((2,  3)\) e
  \((4,  1)\)
\end{question}

%------------------------------------------------------------------------------%
\begin{resolution}{Solução}
  \onslide<+->{a) Determinando a curva de nível}
  \begin{align*}
    \onslide<+->{f(x, y)           & = 4}   \\
    \onslide<+->{(x-2)^2 + (y-1)^2 & = 2^2}
  \end{align*}
  \onslide<+->{Equação da circunferência centrada em \((2, 1)\) e raio 2}
\end{resolution}

%------------------------------------------------------------------------------%
\begin{resolution}{Solução}
  b) Calculando as derivadas parciais de $f$

  \pause
  \[
    \D{f}{x}
    \pause
    \;=\; \D{}{x}\left[(x-2)^2 + (y-1)^2\right]
    \pause
    \;=\; 2(x-2)\D{}{x}(x-2)
    \pause
    \;=\; 2(x-2)
  \]

  \pause
  \[
    \D{f}{y}
    \pause
    \;=\; \D{}{y}\left[(x-2)^2 + (y-1)^2\right]
    \pause
    \;=\; 2(y-1)\D{}{y}(y-1)
    \pause
    \;=\; 2(y-1)
  \]

  \pause
  Gradiente de $f$
  \pause
  \[
    \nabla f(x, y)
    \pause
    \;=\; \nabla\left[(x-2)^2 + (y-1)^2\right]
    \pause
    \;=\; \begin{pmatrix}
      2(x-2) \\ 2(y-1)
    \end{pmatrix}
  \]
\end{resolution}

%------------------------------------------------------------------------------%
\begin{resolution}{Solução}
  Calculando o gradiente nos pontos
  \((0,  1)\) e
  \((2, -1)\)

  \pause
  \[
    \nabla f(0, 1)
    \pause
    \;=\; \begin{pmatrix}
      2(x-2) \\ 2(y-1)
    \end{pmatrix} \evalat{(0, 1)}{}
    \pause
    \;=\; \begin{pmatrix}
    2(0-2) \\ 2(1-1)
    \end{pmatrix}
    \pause
    \;=\; \begin{pmatrix}
    -4 \\ 0
    \end{pmatrix}
  \]

  \pause
  \[
    \nabla f(2, -1)
    \pause
    \;=\; \begin{pmatrix}
      2(x-2) \\ 2(y-1)
    \end{pmatrix} \evalat{(2, -1)}{}
    \pause
    \;=\; \begin{pmatrix}
      2(2-2) \\ 2(-1-1)
    \end{pmatrix}
    \pause
    \;=\; \begin{pmatrix}
      0 \\ -4
    \end{pmatrix}
  \]
\end{resolution}

%------------------------------------------------------------------------------%
\begin{resolution}{Solução}
  Calculando o gradiente nos pontos
  \((2,  3)\) e
  \((4,  1)\)

  \pause
  \[
    \nabla f(2, 3)
    \pause
    \;=\; \begin{pmatrix}
      2(x-2) \\ 2(y-1)
    \end{pmatrix} \evalat{(2, 3)}{}
    \pause
    \;=\; \begin{pmatrix}
      2(2-2) \\ 2(2-1)
    \end{pmatrix}
    \pause
    \;=\; \begin{pmatrix}
      0 \\ 4
    \end{pmatrix}
  \]

  \pause
  \[
    \nabla f(4, 1)
    \pause
    \;=\; \begin{pmatrix}
      2(x-2) \\ 2(y-1)
    \end{pmatrix} \evalat{(4, 1)}{}
    \pause
    \;=\; \begin{pmatrix}
      2(4-2) \\ 2(1-1)
    \end{pmatrix}
    \pause
    \;=\; \begin{pmatrix}
      4 \\ 0
    \end{pmatrix}
  \]
\end{resolution}

%------------------------------------------------------------------------------%
%------------------------------------------------------------------------------%
\begin{question}
  Encontre uma equação para a reta tangente à elipse
  \[
    \frac{x^2}{4} + y^2 = 2
  \]
  no ponto \((-2, 1)\)

  \pause
  \bigskip
  Note que a elipse é uma \\
  curva de nível da função
  \[
    f(x, y) = \frac{x^2}{4} + y^2
  \]

  \only<3>{\addgraphic*{12}{6.5,3}{H-02-Propriedades_gradientes-ex-4}}
\end{question}

%------------------------------------------------------------------------------%
\begin{resolution}{Derivadas parciais de $f$}
  \[
    \D{f}{x}
    \pause
    = \D{}{x}\left(\frac{x^2}{4} + y^2\right)
    \pause
    = \frac{x}{2}
  \]

  \bigskip
  \[
    \pause
    \D{f}{y}
    \pause
    = \D{}{y}\left(\frac{x^2}{4} + y^2\right)
    \pause
    = 2y
  \]
\end{resolution}

%------------------------------------------------------------------------------%
\begin{resolution}{Gradiente}
  Então o gradiente de $f$, no ponto \((-2, 1)\), é
  \pause
  \[
    \nabla f(-2, 1)
    \pause
    \;=\; \Vetor{\frac{x}{2}}{2y}\evalat{(-2, 1)}{}
    \pause
    = \vetor{-1}{2}
  \]
\end{resolution}

%------------------------------------------------------------------------------%
\begin{resolution}{Reta Tangente}
  \onslide<+->{A reta tangente é ortogonal ao vetor gradiente}
  \begin{columns}
  \column{0.5\linewidth}
    \begin{align*}
      \onslide<+->{\nabla f(a, b)  & \cdot \vetor{x-a}{y-b} = 0 \\[3mm]}
      \onslide<+->{\nabla f(-2, 1) & \cdot \vetor{x-(-2)}{y-1}  = 0 \\[3mm]}
      \onslide<+->{\vetor{-1}{2}   & \cdot \vetor{x+2}{y-1}     = 0}
    \end{align*}
  \column{0.5\linewidth}
  \begin{align*}
    \onslide<+->{-(x+2) + 2(y-1) & = 0 \\}
    \onslide<+->{-x +2y -4 & = 0 \\}
    \onslide<+->{-x +2y    & = 4 \\}
    \onslide<+->{ x        & = 2y-4}
  \end{align*}
  \end{columns}
\end{resolution}

%------------------------------------------------------------------------------%
\framedgraphic*[resolution]{Solução}{H-02-Propriedades_gradientes-ex-4_solution}

%------------------------------------------------------------------------------%


%------------------------------------------------------------------------------%
\section{Propriedades Algébricas}

%------------------------------------------------------------------------------%
\begin{frame}{Propriedades Algébricas do Gradiente}
  \begin{enumerate}[<+->]
    \setlength\itemsep{1.5em}
    \item Soma                        \tabto{55mm} \(\nabla(f+g) = \nabla f + \nabla g\)
    \item Diferença                   \tabto{55mm} \(\nabla(f-g) = \nabla f - \nabla g\)
    \item Multiplicação por constante \tabto{55mm} \(\nabla(kf) = k\nabla f\)
    \item Produto                     \tabto{55mm} \(\nabla(fg) = f\nabla g + g\nabla f\)
    \item Quociente                   \tabto{55mm}
      \(\nabla\left(\frac{f}{g}\right) = \frac{g\nabla f - f\nabla g}{g^2}\)
  \end{enumerate}
\end{frame}

%------------------------------------------------------------------------------%
\include{ex/H-02-Propriedades_gradientes-ex5}

%------------------------------------------------------------------------------%
\section{Lista Mínima}

%------------------------------------------------------------------------------%
\begin{frame}{Lista Mínima}
  Cálculo Vol. 2 do Thomas 12\textsuperscript{a} ed. -- Seção 14.5
  \begin{enumerate}
    \item Estudar o texto da seção
    \item Resolver os exercícios: 20, 22, 24, 26, 28, 29
  \end{enumerate}

  \vfill
  Atenção: A prova é baseada no livro, não nas apresentações
\end{frame}

%------------------------------------------------------------------------------%
\end{document}
%------------------------------------------------------------------------------%
